%! Sinusoidal Functions — Unit 6, Lesson 6.6 — Guided Notes
\documentclass[10pt]{article}
\usepackage{precalculus-article}
\usepackage{precalculus-boxes}
\newcommand{\TallMath}[1]{$\displaystyle #1\rule[-1.4em]{0pt}{3.2em}$}

\begin{document}

\pageheader{Unit 6, Lesson 6.6}{Guided Notes}
\namedateperiod

\vspace{0.1in}

%-----------------------------------------------------------------------
% OBJECTIVE
%-----------------------------------------------------------------------
\begin{objectivebox}
\textbf{I will be able to\ldots}
\begin{itemize}
  \item Define a sinusoidal function and explain how cosine is related to sine.
  \item Determine the amplitude, midline, period, and frequency of a sinusoidal
        function from a verbal description or graph.
  \item Describe the symmetry of $y=\sin\theta$ and $y=\cos\theta$.
\end{itemize}
\end{objectivebox}

\vspace{0.08in}

%-----------------------------------------------------------------------
% VOCABULARY
%-----------------------------------------------------------------------
\begin{vocabbox}
\begin{itemize}
  \item \termblanklong{sinusoidal function}{any function formed by additive and
        multiplicative transformations of $f(\theta)=\sin\theta$}
  \item \termblanklong{amplitude}{half the difference between the maximum and
        minimum values of a sinusoidal function}
  \item \termblanklong{midline}{the horizontal line halfway between the maximum
        and minimum; $y = \tfrac{\text{max}+\text{min}}{2}$}
  \item \termblanklong{period}{the smallest positive value $T$ such that
        $f(\theta + T) = f(\theta)$ for all $\theta$}
  \item \termblanklong{frequency}{the reciprocal of the period;
        $\text{frequency} = 1/T$}
  \item \termblanklong{odd function}{a function where $f(-x) = -f(x)$;
        graph has rotational symmetry about the origin}
  \item \termblanklong{even function}{a function where $f(-x) = f(x)$;
        graph has reflective symmetry over the $y$-axis}
\end{itemize}
\end{vocabbox}

\vspace{0.08in}

%-----------------------------------------------------------------------
% HOOK
%-----------------------------------------------------------------------
\begin{hookbox}
\textit{``A recording engineer notices that a sound wave has amplitude 0.3 and
frequency 440 Hz.''}

What does amplitude tell us about the wave?
\writelines{2}

What does frequency tell us?  What is the period of one cycle?
\writelines{2}
\end{hookbox}

\vspace{0.08in}

%-----------------------------------------------------------------------
% CHARACTERISTICS OF SINUSOIDAL FUNCTIONS
%-----------------------------------------------------------------------
\begin{notesbox}{Characteristics of Sinusoidal Functions}

\textbf{Amplitude}
\[
  \text{Amplitude} = \frac{\underline{\hspace{1.5cm}} - \underline{\hspace{1.5cm}}}{2}
  = \text{half the } \underline{\hspace{3cm}}
\]

\textbf{Midline}
\[
  \text{Midline: } y = \frac{\underline{\hspace{1.5cm}} + \underline{\hspace{1.5cm}}}{2}
  = \text{average of } \underline{\hspace{3cm}}
\]

\textbf{Period} \quad The smallest positive \underline{\hspace{3cm}} such that
$f(\theta + k) = f(\theta)$ for all $\theta$.

\textbf{Frequency}
\[
  \text{Frequency} = \frac{\underline{\hspace{1cm}}}{\text{period}}
  = \frac{1}{\underline{\hspace{1.5cm}}}
\]

\medskip
\textbf{For $f(\theta) = \sin\theta$:}

\renewcommand{\arraystretch}{1.8}
\begin{tabular}{@{}llll@{}}
  Amplitude $=$ \blank{1.5cm} &
  Midline $=$ \blank{1.5cm} &
  Period $=$ \blank{1.5cm} &
  Frequency $=$ \blank{2cm}
\end{tabular}

\end{notesbox}

\vspace{0.08in}

%-----------------------------------------------------------------------
% RELATIONSHIP BETWEEN SINE AND COSINE
%-----------------------------------------------------------------------
\begin{notesbox}{Relationship Between Sine and Cosine}

Complete the identity:
\[
  \cos\theta = \sin\!\left(\theta + \underline{\hspace{1.5cm}}\right)
\]

This means cosine is a \underline{\hspace{4cm}} of the sine function.

Therefore cosine is also a \underline{\hspace{4cm}} function.

\end{notesbox}

\vspace{0.08in}

%-----------------------------------------------------------------------
% SYMMETRY
%-----------------------------------------------------------------------
\begin{notesbox}{Symmetry of Sine and Cosine}

\begin{multicols}{2}
\textbf{Sine --- Odd Function}

$\sin(-\theta) = $ \blank{2.5cm}

Graphical meaning: the graph has
\underline{\hspace{3.5cm}}\\
symmetry about the \underline{\hspace{2cm}}.

\columnbreak

\textbf{Cosine --- Even Function}

$\cos(-\theta) = $ \blank{2.5cm}

Graphical meaning: the graph has
\underline{\hspace{3.5cm}}\\
symmetry over the \underline{\hspace{2cm}}.
\end{multicols}

\end{notesbox}

\vspace{0.08in}

%-----------------------------------------------------------------------
% PRACTICE
%-----------------------------------------------------------------------
\begin{practicebox}
\textbf{Practice:} A sinusoidal function has maximum value $5$, minimum value $-1$,
and period $4\pi$.  Find the amplitude, midline, and frequency.

\medskip
Amplitude $= \dfrac{\blank{1.2cm} - (\blank{1.2cm})}{2} = $ \blank{1.5cm}

\medskip
Midline: $y = \dfrac{\blank{1.2cm} + (\blank{1.2cm})}{2} = $ \blank{1.5cm}

\medskip
Frequency $= \dfrac{1}{\blank{2cm}} = $ \blank{2.5cm}
\end{practicebox}

\end{document}
